% Consolidated reply to Reviewer 1bPK. Standalone: compiles on its own.
% Sources for every number: neurips_review/results_table.csv (new runs) and
% figures/flyvis/*_fullsample.tex (full-sample twins of Tabs. 1, 2, Supp. 4, 7).
\documentclass[11pt]{article}
\usepackage[margin=1in]{geometry}
\usepackage{amsmath,amssymb}
\usepackage{booktabs}
\usepackage[colorlinks=true,allcolors=blue]{hyperref}
\setlength{\parskip}{5pt}
\setlength{\parindent}{0pt}
\newcommand{\Rw}{R^2_{\widehat W}}
\newcommand{\Rt}{R^2_{\hat\tau}}
\newcommand{\Rv}{R^2_{\hat V^{\mathrm{rest}}}}
\newcommand{\Mat}[1]{\mathbf{#1}}

\title{}
\date{}
\author{}

\begin{document}
\maketitle
\vspace{-3.5em}

\section*{Reviewer 1bPK}
We thank the reviewer. The weaknesses are well-founded; acting on Q1 and Q2
produced new results.

\paragraph{Q1.}
The three idealizations are coupled and never separated in the paper. We ran two
mismatch experiments; all runs are single-seed, Flyvis-217, $\sigma = 0.05$,
consensus hyperparameters, median $\tau = 19.8$~ms.

\textbf{Model-class match.} A slow adaptation current leaving the hypothesis class
in two ways at once: not first order in the observables, not decomposable into
pairwise messages:
\[
  \tau_i \dot v_i = -v_i + V^{\mathrm{rest}}_i
  + \sum_j W_{ij}\,\mathrm{ReLU}(v_j) - g_a c_i + I_i ,
  \qquad
  \dot c_i = (v_i - c_i)/\tau_a ,
\]
The GNN is unchanged; the target is the observed $\dot v_i$ including the new
term, $\tau_a = 200$~ms $\gg \tau$, $c_i$ never observed.

\begin{center}
{\footnotesize
\begin{tabular}{lcccc}
\toprule
$g_a$ & GNN $\Rw$ / slope & GNN $\Rt$ & GNN $\Rv$ & Known-ODE $\Rw$ \\
\midrule
$0$   & $0.987$ / $0.982$ & $0.99\,(0.99)\,[0.0]$ & $0.97\,(0.63)\,[5.0]$   & $0.986$ \\
$0.1$ & $0.974$ / $0.983$ & $0.99\,(0.99)\,[0.0]$ & $0.80\,(0.62)\,[5.7]$   & $0.972$ \\
$0.3$ & $0.909$ / $0.960$ & $0.98\,(0.98)\,[0.0]$ & $0.54\,(-0.35)\,[27.7]$ & $0.923$ \\
\bottomrule
\end{tabular}}

{\footnotesize $\tau$ / $V^{\mathrm{rest}}$ cells read inlier (full-sample)
[excluded \%].}
\end{center}

Degradation is ordered: $V^{\mathrm{rest}}$ first, $W$ next, $\tau$ last; the
oracle matches. A slow per-neuron offset is what $\widehat V^{\mathrm{rest}}$
absorbs and biases $\widehat W$; $\tau$ is read off the local slope of $f_\theta$.
The priors called class-narrowing are not load-bearing: $\mu_0 = 0$ in every run,
and $\mu_1 = 0$ leaves $\Rw$ at $0.988$ against $0.987$. No sign convention is
imposed: $\widehat W_{ij} \in \mathbb{R}$ carries the sign and $g_\phi^2$ only the
magnitude. Conductance-like coupling needs a
different loss and architecture and is deferred; the adaptation current tests the
same two assumptions. Both mismatches were also applied jointly, $g_a = 0.3$ at
$h = 2$~ms. For Known-ODE their costs in $\Rw$ are near-additive: from $0.955$
matched, adaptation alone costs $0.055$ and the fine step $0.120$; jointly they
cost $0.202$, giving $\Rw = 0.753$ against the $0.781$ a sum would predict. The
GNN pair is still training.

\textbf{Temporal match.} With the model always predicting at $20$~ms:
(i)~generator integrated below the model step, observed at $20$~ms;
(ii)~generator at $2$~ms, state observed every $20$ to $200$~ms.

\begin{center}
{\footnotesize
\begin{tabular}{cccccc}
\toprule
$\Delta t$ generated & $\Delta t$ observed & GNN $\Rw$ / slope & GNN roll.\ $r$ & Known ODE $\Rw$ / slope & Known ODE roll.\ $r$ \\
\midrule
$20$~ms & $20$~ms  & $0.970$ / $0.965$ & $0.999$ & $0.955$ / $1.025$ & $0.999$ \\
$10$~ms & $20$~ms  & $0.864$ / $0.796$ & $0.981$ & $0.873$ / $0.854$ & $0.981$ \\
$4$~ms  & $20$~ms  & $0.822$ / $0.722$ & $0.978$ & $0.844$ / $0.826$ & $0.976$ \\
$2$~ms  & $20$~ms  & $0.812$ / $0.712$ & $0.977$ & $0.835$ / $0.819$ & $0.974$ \\
\midrule
$2$~ms  & $40$~ms  & $0.720$ / $0.717$ & $0.981$ & $0.711$ / $0.723$ & $0.983$ \\
$2$~ms  & $100$~ms & $0.315$ / $0.603$ & $0.980$ & $0.449$ / $0.581$ & $0.986$ \\
$2$~ms  & $200$~ms & (running)         &         & $0.287$ / $0.490$ & $0.980$ \\
\bottomrule
\end{tabular}}
\end{center}

\emph{Fine generation.} $20/20$~ms is the matched case, where the generator
\emph{is} the model's Euler map. A finer step adds the truncation bias
$(\Delta t/2)\,\ddot v$, neuron-specific rather than a global rescaling, costing
$0.16$ $\Rw$.

\emph{Coarse observation} is the real limitation: $\Rw$ falls to $0.287$ at
$10\tau$, reproducing the published $1/5$-frames row ($0.30$ at $100$~ms), while
rollout $r$ stays above $0.97$ throughout, so prediction cannot detect the
violation. Limitation~2 reads ``inference and simulation share $\Delta t =
20$~ms, so the piecewise-linear update is exact rather than approximate''. It
names the idealization but gives no tolerance; we will add the one this sweep
measures, $\Delta t_{\mathrm{obs}} \lesssim \tau$, beyond which $\Mat{W}$ is not
recovered.

\textbf{What we revise.} Known-ODE degrades in step with the GNN, so the cost is in the inverse problem,
not the hypothesis class: the GNN stays within $0.023$ $\Rw$ of it. We retain
``on par with an oracle'' as a claim about the GNN, and drop ``without assuming
the form of the dynamical equation'': the scalar nonlinearities are learned, but
first-order dynamics and additive pairwise aggregation over the supplied graph
are assumed.

\textbf{Simulator-only.} Conceded: no dataset pairs synapse-resolution
connectivity with dense activity in the same animal. We will retitle to
``\dots from simulated neural activity''.

\paragraph{Q2.}
Computed for all four tables (Tabs.~1, 2, Supp.~Tabs.~4 and~7); the twins replace
the submitted ones. Cells read inlier (full-sample) [excluded \%].

\begin{center}
{\footnotesize
\begin{tabular}{llcc}
\toprule
model & $\sigma$ & $\Rt$ in (full) [out.\ \%] & $\Rv$ in (full) [out.\ \%] \\
\midrule
Known ODE & $0$    & $1.00\,(0.09)\,[1.9]$   & $0.97\,(0.86)\,[10.2]$ \\
          & $0.05$ & $1.00\,(1.00)\,[0.0]$   & $0.99\,(0.95)\,[4.7]$ \\
          & $0.5$  & $1.00\,(1.00)\,[0.0]$   & $1.00\,(1.00)\,[0.0]$ \\
\midrule
GNN       & $0$    & $0.94\,(-11.80)\,[3.1]$ & $0.76\,(-0.39)\,[13.7]$ \\
          & $0.05$ & $0.99\,(0.99)\,[0.0]$   & $0.93\,(0.76)\,[3.3]$ \\
          & $0.5$  & $1.00\,(1.00)\,[0.0]$   & $0.98\,(0.97)\,[0.0]$ \\
\bottomrule
\end{tabular}}
\end{center}

Where exclusion is small the two agree. Worst rows of Tab.~2 and Supp.~Tab.~4:

\begin{center}
{\footnotesize
\begin{tabular}{llcc}
\toprule
condition & table & $\Rt$ in (full) & $\Rv$ in (full) [out.\ \%] \\
\midrule
$1/5$ frames                & Supp.~4 & $0.83\,(0.35)$  & $0.49\,(-5.09)\,[56.3]$ \\
meas.\ noise $\gamma = 0.2$ & Supp.~4 & $0.08\,(-0.04)$ & $0.61\,(-5.63)\,[49.1]$ \\
uncert.\ edges, FlyWire eye & Tab.~2  & $0.95\,(0.38)$  & $0.61\,(-9.90)\,[40.8]$ \\
$10\%$ hidden               & Supp.~4 & $0.94\,(-0.87)$ & $0.62\,(-149.37)\,[21.1]$ \\
\bottomrule
\end{tabular}}
\end{center}

\textbf{What we revise.} $\tau$ survives: full-sample $0.99$ against the oracle's
$1.00$ at $\sigma = 0.05$; noise-free it does not ($-11.80$ against $0.09$).
$V^{\mathrm{rest}}$ does not: $0.76$ against $0.95$ at $\sigma = 0.05$, and
negative in every degraded row. We keep ``on par with an oracle'' for $W$ and
$\tau$ at $\sigma > 0$, drop it for $V^{\mathrm{rest}}$, and report degraded rows
as not recovered. On the criterion: filtering must not be defined by the residual
it excuses, so the excluded set will be fixed a priori, from zero in-degree ($44$
neurons) and neurons never crossing threshold; a parameter is then recovered at
full-sample $R^2 > 0.9$, partial below that, not recovered at $R^2 \leq 0$. The
Known-ODE runs were filtered at $\delta_V = 0.1$, not the stated $0.2$, so their
inlier $V^{\mathrm{rest}}$ is twice as tightly filtered as the GNN's; full-sample
is unaffected and we will recompute.

\paragraph{Q3.}
We restrict the claim to $\sigma > 0$ and present the noise-free row as the
ill-posed regime.

\begin{center}
{\footnotesize
\begin{tabular}{lcccc}
\toprule
& GNN $\sigma{=}0$ & Known ODE $\sigma{=}0$ & GNN $\sigma{=}0.05$ & Known ODE $\sigma{=}0.05$ \\
\midrule
$\Rw$               & $0.89$  & $0.96$ & $0.99$ & $0.99$ \\
$\Rv$ (full-sample) & $-0.39$ & $0.86$ & $0.76$ & $0.95$ \\
\bottomrule
\end{tabular}}
\end{center}

On the premise, $\sigma$ is not an experimenter knob but intrinsic stochasticity:
release is probabilistic, channels are noisy, membrane potentials fluctuate. We
have not calibrated its level, so the question is where real circuits sit on the
$\sigma$ axis, not whether $\sigma = 0$.

\paragraph{Q4.}
We have no such baseline and accept the reframing. No random or Bayesian search
was run at matched compute, so the claim that it yields better configurations is
unsupported; it also optimises $\Rw$ against ground truth, unavailable in the
intended application. It becomes a methods note rather than contribution~(3),
with the logs replaced by a curated summary and the self-assessments removed. We
retain only the outcome: one configuration used unchanged everywhere
(Supp.~Tab.~6).

\paragraph{Q5.}
\textbf{Timing.} The \textbf{Allier et al.\ [2026]} preprint appeared February
2026, three months before the deadline: prior work, not concurrent.

\textbf{Different system.} Random recurrent assemblies, no anatomy. Random
weights carry no columnar degeneracy, hence $\Rw = 1.00$ there against $0.89$
here at $\sigma = 0$.

\textbf{Different formulation.} Additively separable there,
$\widehat{\dot v}_i = \phi^*(\Mat{a}_i, v_i)
 + \Omega^*_i(t)\sum_j W_{ij}\,\psi^*(\Mat{a}_i,\Mat{a}_j,v_j)$, nested here,
$\widehat{\dot v}_i = f_\theta(v_i, \Mat{a}_i, m_i, I_i(t))$ with
$m_i = \sum_j \widehat W_{ij}\, g_\phi(v_j,\Mat{a}_j)^2$. Since $f_\theta$ is not
additive in $(v_i, m_i)$, our class is more expressive in how self-state and
drive combine. The delta is not the decomposition but the columnar
identifiability analysis, the degradation battery, the FlyWire-scale result and
the edge ablation; no head-to-head was run.

% Reply to Reviewer 1bPK, Q6 / W5b (Shiu et al. 2024; BrainTrace 2026).
% Facts checked against both papers:
%  - Shiu: LIF on FlyWire, per-edge weights from connectome, sign from
%    neurotransmitter prediction; ONE free parameter W_syn, calibrated so that
%    100 Hz sugar-GRN drive gives ~80% of maximal MN9 firing. 91% of 164
%    testable predictions consistent (84% excluding the split-GAL4 screen).
%  - BrainTrace: Nat Commun (2026) 17:1745. Uses Shiu's LIF dynamics
%    (W_syn = 0.275 mV, T_mbr = 20 ms), >125k neurons, 50M synapses, W from
%    FlyWire, and TRAINS the synaptic weights on 17 min of neuropil calcium.
%    Supervision at 1.2 Hz (833 ms per sample) vs T_mbr = 20 ms, i.e. ~42 tau;
%    loss is at neuropil level, not per neuron; no ground-truth W.

\paragraph{Q6.}
\textbf{Shiu et al.\ [2024]} is a forward model, not an activity fit: weights from
the connectome, sign from neurotransmitter prediction, one global scalar
calibrated once. It is also a counterexample to our own claim that large-scale
simulations have struggled to be predictive: their whole-brain model matches
$91\%$ of $164$ tested predictions. We will correct that sentence. Forward
simulation from a connectome is predictive at brain scale; the inverse direction
remains open.

\textbf{BrainTrace [2026]} is the activity-fitting descendant, and poses the same
inverse problem as our Known-ODE control: the form is fixed (the Shiu LIF), the
connectome supplies the support of $\Mat{W}$, weights optimised so simulated
activity matches recordings. That is the best case in our study, and Appx.~C
shows $\Mat{W}$ is already underdetermined there.

It is also harder: the loss is on \emph{neuropil-level} rates, so each residual
constrains a sum over thousands of neurons, not one cell. Supervision is at $1.2$~Hz,
one sample per $833$~ms against $T_{\mathrm{mbr}} = 20$~ms, i.e.\ $42\tau$: the state relaxes many times between constraints. Our sweep reaches
$10\tau$, where the oracle already falls to $\Rw = 0.287$. With no ground-truth
$\Mat{W}$ for a real fly, fit quality is the only signal, and Appx.~C shows it
cannot distinguish a recovered connectome from a dynamically equivalent one: such
weights are one member of an equivalence class, not recovered connectivity.

Sect.~2 will separate forward models (Shiu) from fits to activity (BrainTrace,
Mi, Pospisil, Lappalainen).

\paragraph{W4.}
Agreed, and it will not stay in Limitations. Here $\gamma$ is the standard
deviation of Gaussian noise added to the voltages after simulation,
$v_i^{\mathrm{obs}} = v_i + \gamma\,\epsilon_i$: unlike $\sigma$ it corrupts the
recording, not the trajectory. At $\gamma = 0.1$, $\Rw = 0.63$ (GNN) and $0.69$
(Known-ODE); at $\gamma = 0.2$, $0.38$ and $0.42$. The oracle degrades as much as
the GNN, so the failure is in the inverse problem rather than in our
architecture, but that does not make the method usable at these levels. The
contributions, not only the Limitations, will say the method is demonstrated
under process noise, not at realistic measurement noise.

\paragraph{W7 and W8.}
Both abstract claims are overstated and will be softened. For W7 the evidence is
intrinsic noise in silico and optogenetics is an untested analogy, so the
abstract will claim only that noise and perturbation help in simulation. For W8
the SIREN rollout is scored on training-frame indices and the stimulus is largely
present in the photoreceptor activities, so ``recovers unknown visual stimuli''
comes out of the abstract. The experiment stays in the appendix, restated as
partial stimulus recovery scored on training-frame indices.

\paragraph{W9.}
\emph{$N$.} $13{,}741$ neurons, $13{,}697$ with at least one incoming edge, $44$
with none: Appx.~C counts the former, so the discrepancy is neurons with no
\emph{input}, not no output. $\sum_i \max(0, d_i - 45) = 115{,}223$ matches
Appx.~C exactly.

\emph{$\pm 0.00$.} Rounding, not under-dispersion. At four decimals the GNN folds
at $\sigma = 0.05$ span $0.9829$ to $0.9886$ (SD $0.0021$), against SD $0.0002$
for Known-ODE, whose fit is least squares in the true parameters. A
byte-identical rerun differs by $0.0014$, so the GNN spread sits near the
non-deterministic GPU floor. We will report three decimals, so these read
$\pm 0.002$. The $0.90 \pm 0.21$ entry is one fold: noise-free GNN rollout $r$ is
$1.000, 1.000, 0.999, 0.482, 0.998$, which the degeneracy predicts, and why we
restrict the claim to $\sigma > 0$.

\end{document}
